\section{Experimental Methodology}
\label{sec:methodology}

\subsection{Evaluation Goals}

The experimental evaluation was designed to answer four related questions.

First, the study evaluates whether relation-aware metadata retrieval improves
Text-to-Spark-SQL correctness compared with providing the complete metadata
catalog to the language model.

Second, it measures whether metadata selection reduces the amount of prompt
context and the corresponding model-evaluation cost.

Third, it studies the behaviour of the retrieval layer when the metadata
catalog grows and contains semantically related distractor sources.

Fourth, it evaluates the Spark execution layer under increasing physical data
volume.

These goals deliberately separate two different notions of scalability:
\emph{catalog scale}, which concerns the number of metadata sources available
to the retrieval system, and \emph{data scale}, which concerns the number of
physical records processed by Spark.


\subsection{Development Benchmark}

An initial benchmark of ten analytical intents was used during development.
The questions cover single-source aggregation, filtering, geographic joins,
cross-taxi comparisons, temporal joins with weather data, and multi-source
queries involving both geographic and weather information.

Each benchmark item contains a gold specification including the metadata
required to answer the question. Depending on the analytical intent, this can
include required datasets, physical columns, relationships, and semantic
rules. The benchmark also contains a gold Spark SQL query and the expected
result.

The initial benchmark was used while implementing and correcting the
retrieval logic. Consequently, after this tuning activity it was explicitly
treated as a \emph{development set} and was not used as the main evidence for
generalisation performance.

This distinction is important because retrieval behaviour had already been
observed on these questions while the architecture was still being modified.


\subsection{Architecture Freeze}

Before the final held-out evaluation, the metadata-aware architecture was
frozen.

The frozen implementation includes:

\begin{itemize}

    \item metadata catalog access;
    \item metadata document construction;
    \item local embedding support;
    \item dense and lexical retrieval;
    \item semantic-concept expansion;
    \item dataset grounding;
    \item relation-aware expansion;
    \item semantic-rule expansion;
    \item SQL-safe context rendering;
    \item SQL validation;
    \item one-shot repair.

\end{itemize}

After this freeze, the retrieval architecture was not tuned using the final
held-out results. Git tags were used to preserve the exact implementation
state associated with each evaluation phase.


\subsection{Frozen Held-Out Benchmark}

A separate held-out benchmark was constructed after the development
architecture had been frozen.

The held-out benchmark contains 20 natural-language questions, with two new
paraphrases for each of the ten analytical intents represented in the
development set. The wording of these questions was not used to tune the
retriever before the benchmark was frozen.

The underlying gold analytical specification for each intent remains fixed:
the expected datasets, columns, relationships, semantic rules, output schema,
gold SQL, and gold result correspond to the same intended computation.
However, the natural-language formulation is new.

The held-out benchmark itself was committed and tagged before the first final
A/B/C system evaluation. This provides a clear separation between benchmark
definition and result observation.


\subsection{Experimental Configurations}

The main Text-to-Spark-SQL experiment compares three configurations.

\paragraph{Configuration A: Full Catalog.}

Configuration A is the baseline. The language model receives the complete
rendered metadata catalog for every question. No metadata retrieval is used to
restrict the catalog before generation.

This baseline represents the straightforward strategy of exposing all
available metadata to the model. It also preserves internal semantic catalog
information that is useful for describing the data but is not necessarily a
physical SQL identifier.

The generated SQL is validated before execution, but no repair attempt is
performed.

\paragraph{Configuration B: Relation-Aware Retrieval.}

Configuration B uses the proposed metadata retrieval architecture. Dense and
lexical evidence identify candidate metadata, after which semantic and
relationship expansion add structurally required entities.

Only the selected metadata is rendered into the SQL-safe context. The same
language model and generation parameters used by Configuration A are then used
to generate Spark SQL.

The generated SQL is validated before execution, but no repair attempt is
performed.

\paragraph{Configuration C: Retrieval with Validation and Repair.}

Configuration C uses the same relation-aware retrieval and SQL-safe context as
Configuration B.

The difference is that, when the initially generated SQL fails validation, the
system can perform one repair attempt using the validator feedback. If the
query is already valid, no repair is triggered.

This design isolates the incremental contribution of one-shot repair from the
contribution of metadata retrieval.


\subsection{Controlled Generation Configuration}

All A/B/C Text-to-Spark-SQL experiments use the same local
\texttt{qwen2.5-coder:3b} model through Ollama.

The principal generation parameters are fixed to:

\begin{itemize}

    \item temperature \(=0\);
    \item context size \(=4096\) tokens.

\end{itemize}

Using the same model and deterministic temperature across configurations
reduces the number of variables changed by the A/B comparison. The primary
difference between A and B is therefore the metadata context supplied to the
model.


\subsection{Metadata Retrieval Metrics}

Metadata retrieval is evaluated independently from SQL generation.

For each held-out query, the selected metadata entities are compared with the
gold required-metadata set.

The evaluation uses:

\begin{itemize}

    \item \textbf{recall}, measuring the fraction of required metadata entities
    successfully retrieved;

    \item \textbf{precision}, measuring the fraction of retrieved metadata
    entities that belong to the gold required set;

    \item \textbf{F1}, combining precision and recall;

    \item \textbf{perfect-recall queries}, counting questions for which all
    required metadata entities are present;

    \item \textbf{context size}, measured in characters;

    \item \textbf{context reduction}, measured relative to the complete
    catalog context;

    \item \textbf{retrieval latency}.

\end{itemize}

Gold metadata can contain physical datasets, columns, relationships, and
semantic rules. As a consequence, retrieval evaluation measures structural
completeness rather than only nearest-neighbour relevance.


\subsection{Text-to-Spark-SQL Metrics}

Generated SQL is evaluated at multiple levels.

\textbf{Validation rate} measures how often the generated query passes the
parser, read-only constraints, physical schema analysis, and expected-output
checks.

\textbf{Execution rate} measures how often a generated query can be executed
successfully using Spark.

The principal correctness metric is \textbf{result accuracy}. A generated
query is considered correct when its executed result matches the expected
gold result under the benchmark comparison procedure.

Result accuracy is preferred over exact SQL-string matching because different
SQL formulations can express the same computation. Conversely, a query can be
syntactically valid and executable while still returning the wrong analytical
result. Validation success is therefore not treated as a substitute for
semantic correctness.


\subsection{Prompt-Efficiency Metrics}

Prompt efficiency is evaluated by measuring the metadata-dependent input
provided to the local model.

The held-out A/B experiment records mean prompt-token count for both the
complete-catalog baseline and the relation-aware configuration.

The relative token reduction is calculated as:

\[
R_{\mathrm{tokens}}
=
1 -
\frac{T_B}{T_A},
\]

where \(T_A\) is the mean prompt-token count for the full-catalog baseline and
\(T_B\) is the mean count for relation-aware prompting.

Retrieval latency is measured separately so that the cost of metadata
selection is not hidden inside the LLM timing.


\subsection{Controlled Latency Protocol}

The original held-out execution revealed that model prompt-evaluation timing
can be strongly influenced by runtime cache effects. In particular, repeated
full-catalog prompts share a large common prefix, which can make later
requests appear artificially inexpensive.

For this reason, intrinsic A/B prompt cost is evaluated using a separate
controlled latency benchmark rather than relying only on the sequential
held-out timings.

Three representative held-out questions with perfect metadata recall were
selected:

\begin{itemize}

    \item one simple single-source analytical request;

    \item one taxi--weather temporal query;

    \item one three-source query involving taxi, geography, and weather.

\end{itemize}

Each A/B condition is executed twice in the controlled phase. A unique
prefix-isolation header is inserted before the otherwise frozen prompt
structure so that reuse of an identical long prefix is reduced.

The benchmark records prompt-token count, prompt-evaluation time, generation
time, API duration, and wall-clock time.

A separate natural-cache phase executes the exact prompts without the
isolation header. These measurements are retained to illustrate the practical
effect of prefix reuse, but the prefix-isolated measurements are used as the
primary latency comparison.

Model cold-load time is measured separately and is not mixed with the
warm-model prompt-evaluation comparison.


\subsection{Repair Evaluation}

Repair is evaluated only when Configuration C encounters an initially invalid
query.

The benchmark records:

\begin{itemize}

    \item the number of repair attempts;

    \item the number of repaired queries that pass validation;

    \item whether a repaired query improves final result accuracy.

\end{itemize}

A repair is not attempted for a query that is already valid. Consequently,
the repair mechanism cannot correct a semantically incorrect query that passes
validation.

Likewise, repair cannot reliably recover from missing metadata if the
retrieval stage failed to supply the physical schema or relationship required
for the intended computation.


\subsection{Catalog-Scale Experiment}

A separate experiment evaluates metadata retrieval as the catalog grows.

Three catalog sizes are considered:

\begin{itemize}

    \item 4 sources: the original physical data lake;

    \item 8 sources: the original catalog plus 4 metadata-only hard-negative
    sources;

    \item 16 sources: the original catalog plus 12 metadata-only
    hard-negative sources.

\end{itemize}

The additional sources are deliberately semantically related to the original
urban-mobility, geographic, and weather domain. They include concepts such as
alternative trip data, borough information, daily weather, transportation,
traffic, and environmental observations.

These sources are metadata-only distractors: no additional physical Spark data
is executed in this experiment. The purpose is to increase catalog ambiguity
while keeping the analytical gold answers unchanged.

The distractor specification and the evaluation harness were frozen before
the first catalog-scale run. No retrieval tuning was performed after observing
the 4/8/16-source results.

For each scale, a new metadata index is built and the same 20 frozen held-out
questions are evaluated with the same dense Top-\(5\) retrieval setting.

The experiment records:

\begin{itemize}

    \item number of datasets, columns, and metadata documents;

    \item complete-catalog context size;

    \item mean retrieved-context size;

    \item perfect-recall query count;

    \item macro precision, recall, and F1;

    \item mean context reduction;

    \item retrieval latency;

    \item metadata-index construction cost.

\end{itemize}


\subsection{Catalog-Scale Error Analysis}

After the catalog-scale results had been frozen, failures were analysed by
comparing the missing and extra metadata entities at 4, 8, and 16 sources.

This analysis is descriptive rather than corrective. The retrieval logic was
not modified in response to the observed failures.

The purpose of this analysis is to identify whether a loss of recall is caused
by general catalog growth or by competition from particular semantically
similar hard-negative metadata.


\subsection{Data-Volume Scaling Experiment}

Catalog growth and physical data growth are evaluated separately.

The data-volume benchmark uses the curated Yellow Taxi dataset and four row
counts:

\begin{itemize}

    \item 100,000 rows;
    \item 500,000 rows;
    \item 1,000,000 rows;
    \item the complete 2,963,713-row curated Yellow Taxi dataset.

\end{itemize}

Each scale is materialized outside the timed query execution as a two-file
Parquet dataset. The number of partitions is therefore held constant across
the measured scales.

Three Spark SQL workloads are evaluated:

\begin{enumerate}

    \item a scan with aggregate statistics;

    \item grouping by pickup zone;

    \item a temporal join between Yellow Taxi pickups and hourly weather.

\end{enumerate}

Each workload is executed three times for every scale. Scale order is
alternated between repetitions to reduce systematic ordering effects, and no
DataFrame is intentionally cached by the benchmark.

The principal reported timing is the median of the three repetitions. Median
runtime is used because short local Spark workloads are sensitive to
one-time JVM, filesystem, and execution warm-up effects.

The benchmark also records physical Parquet size and runtime relative to the
100,000-row baseline.


\subsection{Interpretation Constraints}

The experiments are intended to evaluate the implemented prototype, not to
establish universal performance claims for Spark, Qwen, Qdrant, or the
underlying embedding model.

In particular, language-model latency is specific to the CPU-only Codespaces
environment, and the Spark data-volume experiment is limited to approximately
three million Yellow Taxi records on two local worker threads.

Similarly, the catalog-scale distractors are a controlled set of
domain-related hard negatives rather than a complete simulation of an
enterprise metadata catalog.

The results are therefore interpreted as evidence about the behaviour and
trade-offs of the proposed architecture under the defined experimental
conditions.
